\mxnewpage\section*{Correction EV2-B}\label{EV2-B-CORR}
\textbf{QCM :} 1 A, 2 C, 3 D, 4 B, 5 A, 6 B, 7 C, 8 C, 9 A, 10 D, 11 D, 12 C. Chaque réponse correcte : \num{0,5} point ; total 6.
\paragraph{Vérifications des automatismes.} Priorités ; somme de fractions ; soustraction des exposants ; coefficient multiplicateur \num{1,2} ; division par \num{0,8} ; produit \num{1,25} fois \num{0,8} ; résolution affine ; accroissement vertical divisé par l’horizontal ; moyenne ; angle divisé par 360 ; complément à 1 ; substitution entre parenthèses.
\subsection*{Exercice 1 -- Indépendance et répétition}
\begin{enumerate}
\item $\dfrac{1}{6}$ et $\dfrac{1}{4}+\dfrac{2}{3}-\dfrac{1}{6}=\dfrac{3}{4}$. \textit{Barème : 1 point. Résultat et justification élémentaire attendus ; demi-point possible pour démarche correcte inachevée.}
\item Trois chemins, chacun $(\dfrac{1}{4})^2(\dfrac{3}{4})$ : total $\dfrac{9}{64}$. \textit{Barème : 2 points. Méthode ou justification : 1 ; calcul et conclusion : 1.}
\item Complément de trois échecs : $1-(\dfrac{3}{4})^3=\dfrac{37}{64}$. \textit{Barème : 2 points. Méthode ou justification : 1 ; calcul et conclusion : 1.}
\end{enumerate}
\subsection*{Exercice 2 -- Rationnel et symétrie}
\begin{enumerate}
\item $f=x+\dfrac{4}{x}$, impaire. \textit{Barème : 1 point. Résultat et justification élémentaire attendus ; demi-point possible pour démarche correcte inachevée.}
\item $f\prime =\dfrac{x^2-4}{x^2}$ ; décroît jusqu’à 2 puis croît. \textit{Barème : 2 points. Méthode ou justification : 1 ; calcul et conclusion : 1.}
\item Croît sur $\intof{-\infty}{-2}$, décroît sur $\intfo{-2}{0}$. Sur $\intff{1}{2}$, f décroît : contre-exemple à la croissance globale. \textit{Barème : 2 points. Méthode ou justification : 1 ; calcul et conclusion : 1.}
\end{enumerate}
\subsection*{Exercice 3 -- Distance et paramètre}
\begin{enumerate}
\item Normale (3;4) ; $12-12=0$. \textit{Barème : 1 point. Résultat et justification élémentaire attendus ; demi-point possible pour démarche correcte inachevée.}
\item $\lambda =-\dfrac{12}{25}$, donc $H\pt{\dfrac{36}{25};\dfrac{48}{25}}$. \textit{Barème : 1 point. Résultat et justification élémentaire attendus ; demi-point possible pour démarche correcte inachevée.}
\item Distance $\dfrac{12}{5}$ ; $x^2+y^2=\dfrac{144}{25}$. \textit{Barème : 2 points. Méthode ou justification : 1 ; calcul et conclusion : 1.}
\end{enumerate}
\paragraph{Étalonnage.} Durée cible : 20–25 min pour les automatismes, 80–85 min pour les exercices, 10–15 min de relecture. Ces durées sont des estimations éditoriales à mesurer en classe. Accepter toute méthode correcte ; ne pas sanctionner deux fois une erreur numérique propagée avec cohérence. Le barème proposé ne vaut pas décision d’harmonisation de l’équipe.
